% Матрица рисков ИИ-методиста
\begin{figure}[htbp]
  \centering
  \begin{tikzpicture}
    \begin{axis}[
      width=0.88\textwidth,
      height=0.66\textwidth,
      xmin=0.5, xmax=3.5,
      ymin=0.5, ymax=3.5,
      xtick={1,2,3},
      ytick={1,2,3},
      xticklabels={Низкая, Средняя, Высокая},
      yticklabels={Среднее, Высокое, Критическое},
      xlabel={Вероятность},
      ylabel={Влияние},
      label style={font=\small},
      tick label style={font=\small},
      grid=major,
      grid style={thin, gray!50},
      axis on top,
      clip=false,
    ]
    % Цветовые зоны фона
    % Зелёная зона (низкая вер. + среднее влияние)
    \fill[green!15] (axis cs:0.5,0.5) rectangle (axis cs:1.5,1.5);
    % Жёлтая зона
    \fill[yellow!15] (axis cs:1.5,0.5) rectangle (axis cs:2.5,1.5);
    \fill[yellow!15] (axis cs:0.5,1.5) rectangle (axis cs:1.5,2.5);
    \fill[yellow!15] (axis cs:1.5,1.5) rectangle (axis cs:2.5,2.5);
    % Оранжевая зона
    \fill[orange!15] (axis cs:2.5,0.5) rectangle (axis cs:3.5,1.5);
    \fill[orange!15] (axis cs:2.5,1.5) rectangle (axis cs:3.5,2.5);
    \fill[orange!15] (axis cs:0.5,2.5) rectangle (axis cs:1.5,3.5);
    % Красная зона
    \fill[red!15] (axis cs:1.5,2.5) rectangle (axis cs:2.5,3.5);
    \fill[red!15] (axis cs:2.5,2.5) rectangle (axis cs:3.5,3.5);

    % Риски из матрицы презентации. Небольшой сдвиг по X сохраняет читаемость
    % подписей при одинаковой средней вероятности.
    \addplot[only marks, mark=*, mark size=5pt, red!70] coordinates {(1.82, 3.00)};
    \node[font=\scriptsize, align=center, text width=2.0cm, anchor=south, yshift=4pt]
      at (axis cs:1.82,3.00) {152-ФЗ};

    \addplot[only marks, mark=*, mark size=5pt, orange!80!black] coordinates {(2.18, 2.80)};
    \node[font=\scriptsize, align=center, text width=2.2cm, anchor=south west, xshift=2pt, yshift=2pt]
      at (axis cs:2.18,2.80) {Финансовый\\риск};

    \addplot[only marks, mark=*, mark size=5pt, blue!70] coordinates {(1.82, 2.30)};
    \node[font=\scriptsize, align=center, text width=2.0cm, anchor=north east, xshift=-2pt, yshift=-3pt]
      at (axis cs:1.82,2.30) {ФГОС};

    \addplot[only marks, mark=*, mark size=5pt, purple!70] coordinates {(2.20, 2.15)};
    \node[font=\scriptsize, align=center, text width=2.4cm, anchor=south west, xshift=2pt, yshift=2pt]
      at (axis cs:2.20,2.15) {Авторские\\права};

    \addplot[only marks, mark=*, mark size=5pt, teal!70!black] coordinates {(1.84, 1.55)};
    \node[font=\scriptsize, align=center, text width=2.4cm, anchor=north east, xshift=-2pt, yshift=-3pt]
      at (axis cs:1.84,1.55) {Конкурентный\\риск};

    \addplot[only marks, mark=*, mark size=5pt, green!50!black] coordinates {(2.18, 1.20)};
    \node[font=\scriptsize, align=center, text width=2.0cm, anchor=south west, xshift=2pt, yshift=2pt]
      at (axis cs:2.18,1.20) {Командный\\риск};

    \end{axis}
  \end{tikzpicture}
  \caption{Матрица рисков ИИ-методиста}
  \label{fig:risk-matrix}
\end{figure}
